% FILE: 06_contribution_to_thesis.tex
% PROJECT: crosswalks
% ROLE: type-law-crosswalk-between-standards-and-execution

\section{Contribution to the Dissertation}
\label{sec:crosswalks-contribution}

\subsection{Contribution to Prediction vs.~Coordination}
\label{sec:crosswalks-contribution-prediction}

The central thesis argument distinguishes between systems that \emph{predict} what a process will
do and systems that \emph{coordinate} what a process lawfully may do. Prediction systems imitate
surface distributions; coordination systems enforce structural constraints. The crosswalks project
contributes to this argument by demonstrating that the translation layer between process formalisms
is itself a site where this distinction matters.

When \textsc{PM4Py} converts \textsc{OCEL} to \textsc{XES} by silently selecting a case notion,
it performs a prediction-style operation: it chooses what looks like a complete event log without
coordinating with the analyst about which object types are being discarded. The resulting
\textsc{XES} log is a prediction of what the analyst wanted, not a coordination with the analyst's
explicit structural intent.

The \texttt{OCEL\_TO\_XES.md} crosswalk replaces this with a coordination contract: the analyst
must declare a \texttt{ProjectionName}, must acknowledge the \texttt{LossReport}, and must
explicitly permit the loss before the projection proceeds. The projection output is not a
prediction of the intended artifact --- it is the artifact the analyst explicitly authorised,
with the authorisation recorded. This is the structural difference between prediction and
coordination at the translation layer.

\subsection{Contribution to Process Evidence}
\label{sec:crosswalks-contribution-evidence}

The crosswalks project establishes that process evidence has a provenance structure that extends
beyond the event log itself. An \textsc{OCEL} log projected to \textsc{XES} without a
\texttt{LossReport} is evidence that has been laundered: its structural ancestry has been erased.
The receipt chain that binds it cryptographically binds a laundered artifact, not the original
multi-object process.

The \texttt{alive\_001\_to\_iso\_compliance.md} crosswalk formalises the evidence integrity
requirement under \textsc{ISO/IEC} 27001: every evidence block must carry a \textsc{BLAKE3}
hash over its serialised payload, state, witness, and epoch, verified by an Ed25519 signature.
The crosswalk connects this cryptographic requirement to the \textsc{ALIVE\_001} gate on
cryptographic binding, grounding the receipt doctrine in an international information security
standard.

\textbf{[AUTHOR THESIS]:} The crosswalks project contributes to process evidence by defining
the \emph{pre-receipt} layer: the admission, projection, and loss-accounting steps that must
complete before an artifact is eligible to enter the receipt chain. Without this layer, a
receipt chain is a chain of laundered evidence. With it, a receipt chain is a chain of
explicitly authorised, loss-acknowledged, structurally grounded artifacts.

\subsection{Contribution to Manufacturing Architecture}
\label{sec:crosswalks-contribution-manufacturing}

The crosswalks project contributes to the manufacturing architecture of the process intelligence
foundry by defining the type-law surface that \texttt{wasm4pm-compat} must enforce. The
manufacturing pipeline in the Chatman Equation framework is:

\[
  \text{raw external data}
  \to \text{admitted compat}
  \to \text{explicit projection with LossPolicy}
  \to \text{LossReport}
  \to \text{receipt-eligible artifact}
\]

The crosswalks project specifies each operator in this pipeline for each formalism pair. The
\texttt{PROCESS\_TREE\_TO\_WFNET.md} crosswalk specifies the lossless operator; the
\texttt{BPMN\_TO\_WFNET.md} specifies the conditionally lossy operator with named refusals;
the \texttt{OCEL\_TO\_XES.md} specifies the inherently lossy operator with mandatory
\texttt{ProjectionName} audit.

The Leemans loop arity enforcement --- \texttt{TypedLoopNode<const ARITY: usize>} with
\texttt{Require<\{ ARITY == 2 \}>: IsTrue} --- is a concrete example of how a manufacturing
architecture can enforce a mathematical theorem (Leemans, 2022, loop definition) at the
type-system level, making it impossible to manufacture a malformed loop node.

\subsection{Contribution to Capital-Flow Infrastructure}
\label{sec:crosswalks-contribution-capital}

\textbf{[AUTHOR THESIS / INTERPRETATION]:} The crosswalks project contributes to capital-flow
infrastructure by providing the type-law foundation for cross-formalism evidence that can
underpin M\&A diligence claims. An M\&A claim based on a conformance score is only admissible
as a board-level claim if the conformance score is grounded in a lawful evidence chain. If the
WF-net used for conformance checking was projected from a \textsc{BPMN} process with silent
\textsc{OR}-gateway approximation, the score is not grounded --- and the claim is not admissible
(source: \texttt{ALIVE\_GATE\_ASSESSMENT.md}, ``What Is Not Authorized'': ``Conformance scores
without a named witness'').

The \textsc{ALIVE\_001} gate-to-\textsc{ISO} crosswalk (\texttt{alive\_001\_to\_iso\_compliance.md})
connects the foundry's checkpoint gates to \textsc{ISO} 9001:2015 (replay soundness), \textsc{ISO/IEC}
23745 (schema completeness), \textsc{ISO/IEC} 27001 (cryptographic binding), and \textsc{ISO/IEC}
12207:2017 (lifecycle state soundness). These are the compliance mappings that an M\&A diligence
package would cite as the regulatory basis for process intelligence claims. The crosswalks project
is the research artefact that grounds those citations.
